\problemname{Dragkamp}
Under en idrottslektion vill läraren dela in alla elever i två lag för att tävla i dragkamp.
För att göra dragkampen intressant bör lagen vara exakt lika stora, samt exakt lika starka.
Läraren har på förhand uppskattat varje elevs styrka till ett tal mellan $1$ och $4$.

Givet hur många elever av varje styrketal som är med på lektionen, kan du berätta hur läraren ska dela upp lagen, om möjligt?

\section*{Indata}
Den första och enda raden innehåller de fyra heltalen $s_1, s_2, s_3, s_4$ ($0 \le s_i \le 10^{7}$), där $s_i$ är antalet elever som har styrkan $i$.

\section*{Utdata}
Om det inte är möjligt att dela in de två lagen, skriv ut en rad med texten ''\texttt{nej}''.

Annars, skriv ut en rad med fyra heltal $a_1, a_2, a_3, a_4$, där $a_i$ är antalet elever med styrkan $i$ som ska vara med i det ena dragkampslaget.

\section*{Poängsättning}
Din lösning kommer att testas på en mängd testfallsgrupper.
För att få poäng för en grupp så måste du klara alla testfall i gruppen.

\noindent
\begin{tabular}{| l | l | p{12cm} |}
  \hline
  \textbf{Grupp} & \textbf{Poäng} & \textbf{Gränser} \\ \hline
  $1$    & $20$      & $s_i \le 300$, en giltig indelning existerar. \\ \hline
  $2$    & $20$      & Ingen elev har styrka $4$, en giltig indelning existerar. \\ \hline
  $3$    & $20$      & En giltig indelning existerar. \\ \hline
  $4$    & $40$      & Inga ytterligare begräsningar. \\ \hline
\end{tabular}

\section*{Förklaring av exempelfall}
TODO
